% Unnumbered geometric matching note. No arbitrary-coefficient hull operation.
\subsection*{Canonical-order stability and a controlled relative full replacement}

This note records three exact pointwise implications for a contracted
$A$-chain. The actual identification of an arbitrary geometric gerbe with
the specified quantum toric model is a separate hypothesis. No global
small object atlas or unbounded-period hereditary count is asserted.

\paragraph{Degree-zero lines and extension of the order action.}
Let $Y\to\operatorname{Spec}S$ be the minimal complete $A_{N-1}$
resolution, and let $D=\sum_{i=1}^{N-1}E_i$ be its reduced exceptional
chain. For $N>2$, the line $\mathcal O_D(-nD)$ has multidegrees
$(n,0,\ldots,0,n)$; for $N=2$ its degree is $2n$.
It has $H^1=0$ for every $n\ge1$, by normalization of the rational tree
and surjective evaluations at its nodes. The square-zero ideal for
$nD\subset(n+1)D$ is $\mathcal O_D(-nD)$, and coherent $H^2$ on $D$
vanishes. The units sequence therefore makes
$\operatorname{Pic}((n+1)D)\to\operatorname{Pic}(nD)$ an isomorphism.
A line bundle with degree zero on every $E_i$ is trivial on $D$ and
hence on the full formal neighborhood, with compatible trivializations
(the same $H^1$ vanishing also lifts units). Thus an inherited line
$\mathcal O(-Z)$ with $Z\cdot E_i=0$ is formally trivial here. Its
pullback to a toric power cover is still trivial.

Let $P$ be an actual locally free order, pulled from the original strong
comparable-cycle model, so that $P(-Z)$ has the proved ordinary global
generating quotient. On this formal block the preceding trivialization
makes this a finite quotient $\mathcal O^d\twoheadrightarrow P$.
Suppose a given punctured twisted anchor has been extended by the exact
quantum toric full functor, producing $A_{\mathrm{can}}$ whose pullback
upstairs is an ordinary full bundle. Each of those generating sections
of $P$ acts on the punctured anchor. Full-Hom extends its action uniquely
to $A_{\mathrm{can}}$. Every linear relation among the generators acts
as zero on the dense puncture and hence everywhere, because the target
endomorphism sheaf is locally free. The maps therefore factor through
an actual $P$-action. Unit and associativity identities extend in the
same way. Equivariance is retained by uniqueness of full-Hom extension.
This proves stability of the new full anchor without requiring an
inclusion from the previous anchor. It is a statement within the
specified full category and does not identify two arbitrary lattices.

\paragraph{Averaging an ordinary source factorization.}
Let $Q$ be a right $P$-module and suppose
\[
 h:Q\longrightarrow\mathcal O^n,\qquad
 k:\mathcal O^n\longrightarrow Q,\qquad kh=z\operatorname{id}_Q
\]
are ordinary linear maps. Let the actual integral scaled separability
tensor be
\[
 e=\sum_i u_i\otimes v_i\in P\otimes P,
 \quad \sum_i a u_i\otimes v_i=\sum_i u_i\otimes v_i a,
 \quad \sum_i u_iv_i=d\,1.
\]
The right action in the second tensor factor is the opposite-algebra
convention. Define
\[
 \Phi(q)=\sum_i h(q u_i)\otimes v_i\in\mathcal O^n\otimes P,
 \qquad \Psi(v\otimes a)=k(v)a.
\]
The centralizing identity gives $\Phi(qa)=\Phi(q)a$, while $\Psi$ is
plainly right $P$-linear. Direct multiplication gives
\[
 \Psi\Phi(q)=\sum_i kh(q u_i)v_i
            =zq\sum_i u_iv_i=dzq.
\]
This formula uses the integral tensor $e=d\sigma$ as a whole. It does
not assume that a chosen decomposition of the generic separability
tensor has separately integral factors.

For the proved full orders we apply this formula separately to the
\emph{global} pencil functions $z=x$ and $z=y$. For each choice,
(D.SN.2) supplies the actual global tensor $e_z$ with product $z^2$,
and the original Noether factorization of $z\operatorname{id}_Q$
pulls back to global maps $h_z,k_z$ through a finite free bundle.
Thus $x^3\operatorname{id}_Q$ and $y^3\operatorname{id}_Q$ each factor
through some $P^{n_z}$ by actual global $P$-linear maps. These maps
restrict to the \emph{whole} formal contracted block and its toric
cover. No map defined only on part of a pencil chart is substituted
into the global full-Hom extension theorem below. The local normalized
tensor (D.SN.3) is not needed for this passage.

Give $A_{\mathrm{can}}$ its left $P$-action and form the torsion-free
quotient of $Q\otimes_P A_{\mathrm{can}}$, followed by its reflexive
hull $F$ on the regular root surface. Tensoring $\Phi,\Psi$ gives
maps to and from $A_{\mathrm{can}}^n$. The first map kills torsion
because its target is locally free. Away from a finite set the
resulting torsion-free module equals its hull. Maps between the
locally free targets extend uniquely across that finite set on the
regular two-dimensional root charts. Consequently
\[
        z^3\operatorname{id}_F:
          F\longrightarrow A_{\mathrm{can}}^n\longrightarrow F
                                                               \qquad(1)
\]
is an actual factorization. Equivariance and all identities extend
by uniqueness. This is geometric reflexive extension over the field;
no coefficient-flat reflexive-hull functor is claimed.

\paragraph{Comparison to the relative full hull.}
Let $\widetilde F$ be the ordinary pullback of $F$ to the toric cover,
and let $\widetilde W$ be the ordinary full extension of its punctured
bundle on the contracted RDP. Suppose $\widetilde F$ is generically
generated along the exceptional chain by copies of the ordinary full
$\widetilde A_{\mathrm{can}}$. This hypothesis holds for the induced
hull above if the actual quotient generating $Q(-Z)$ is retained:
after trivializing $Z$, the maps from copies of $A_{\mathrm{can}}$
generate the tensor quotient, and its hull agrees at every height-one
point.

Each generating map from $\widetilde A_{\mathrm{can}}$ to
$\widetilde F$ on the puncture extends to $\widetilde W$ by full-Hom.
Relations vanish on the puncture. At every exceptional generic point
the generating maps surject onto $\widetilde F$, so the identity
comparison is regular there. It is regular off the exceptional set as
well, and extends over the remaining finite set. Thus
\[
                         \widetilde F\subset\widetilde W. \qquad(2)
\]
Apply (1) separately to the global $x^3$ and $y^3$
factorizations just constructed. Their first maps are defined on the
whole formal resolution, hence on the whole puncture of the contracted
complete RDP. Full-Hom extends them from $\widetilde W$ to their
respective full targets $\widetilde A_{\mathrm{can}}^{n_z}$.
Composing with the unchanged second maps gives multiplication by
$x^3$ and $y^3$, respectively. Thus both inclusions
$x^3\widetilde W\subset\widetilde F$ and
$y^3\widetilde W\subset\widetilde F$ hold globally on the block. Since
$(x,y)\mathcal O=\mathcal O(-C)$, also
$(x^3,y^3)\mathcal O=\mathcal O(-3C)$, so
\[
           \widetilde W(-3C)\subset\widetilde F
                                      \subset\widetilde W. \qquad(3)
\]
Here and below $C$ denotes its actual pullback upstairs. Exact tame
invariants descend this comparison within the specified root category.
The bound is independent of the anchor rank and the chain length.

For an ordinary lattice whose $z\operatorname{id}$ already factors
through ordinary full targets, the same proof gives (3) with $3C$
replaced by $C$. In particular the ordinary canonical full bundles
and their actual blowup pullbacks satisfy this on a zero-$Z$ block.
The statement is a conductor inclusion, not an assertion that the
factorization itself is a generating quotient of $F(-3C)$.

\paragraph{What happens to a prescribed one-step unit map.}
Suppose actual lattices $F_i$ carry a prescribed rational map
$\varphi:F_i^{(p)}\to F_{i+1}$ with forward poles at most $aC$ and
inverse poles at most $bC$. If their relative full replacements satisfy
$W_i(-dC)\subset F_i\subset W_i$ for a fixed integer $d$, then the
same rational map has forward poles at most $(a+pd)C$ and inverse poles
at most $(b+d)C$ between $W_i^{(p)}$ and $W_{i+1}$.
Indeed $W_i^{(p)}(-pdC)\subset F_i^{(p)}$ for the forward map, and
$W_{i+1}(-dC)\subset F_{i+1}$ for the inverse. All maps and their
punctured identification are retained. With (3), the additional losses
are $3pC$ and $3C$.

The old rank-one calculation also shows why the inherited cycle matters.
For $I_a=(u,w^a)$, write $pa=kN+c$, $0\le c<N$. The one-step defect
between $F_a^{(p)}$ and the actual full reflexive target $u^kF_c$ is
\[
       p\min(j,a)-kj-\min(j,c)=p(G_N)_{ja}-(G_N)_{jc}.
\]
It can grow linearly in $N$. No fixed fundamental-cycle bound follows.
In the actual contracted-chain setting of the existing Jacobian proof,
however, (D.LR.8) gives $\min(j,N-j)\le\operatorname{ord}_{E_j}R$
at every vertex, and $R\le\lambda C$. Together with the explicit
bound $d_j\le p\min(j,N-j)$, this gives the required fixed multiple
of the inherited $C$. The contraction must descend to the fixed
regular base model; arbitrary extra blowups are not covered.

\paragraph{Remaining hypotheses and source scope.}
The actual toric gerbe identification, its scalar-band and end-framing
data, and the resulting compatibility with the neighboring core have
not been established by this note. Nor does it construct a small
all-coefficient core chart or control large thickened core overlaps.
Only actual geometric hulls are normalized here; arbitrary Artinian
hulls and their flatness are not inferred. Once objects lie in the
specified full toric category, its previously checked Artinian rigidity
and first-coarse-jet framing apply.

The complete scaled-separability proof, the full existing RDP Hom and
Artinian arguments, and the full Jacobian/linear-resolution proof were
freshly reread. The averaging identities, pointwise extension and pole
calculations are written out above. Earlier exact primary audits of
those inputs are reused; no new external theorem or primary read is
claimed. This is an unnumbered research note and has not been promoted.
